% !Mode:: "TeX:UTF-8"
% parts/abstract.tex

%%%--- "中文摘要" --- %%%%%%%%%%%%%%%%%
% 硕士学位论文中文摘要：约500字左右

\begin{flushleft}
{\zihao{3}\heiti
    \makebox[5em][s]{\textbf{论文题目：}}\textbf{\the\CTitle}\\
    \makebox[5em][s]{\textbf{专业名称：}}\textbf{\the\Cmajor}\\
    \makebox[5em][s]{\textbf{研究生：}}\textbf{\the\StudentAuthor}\hfill \makebox[3.5em][s]{\textbf{签名：}}\thickuline{\quad\quad\quad\quad\quad\quad}\\
    \makebox[5em][s]{\textbf{指导老师：}}\textbf{\the\Csupervisor~~\the\CsupervisorTitle}\hfill \makebox[3.5em][s]{\textbf{签名：}}\thickuline{\quad\quad\quad\quad\quad\quad}\\
}
\end{flushleft}

\begin{cnabstract}
硕士学位论文摘要应说明研究目的、研究方法、成果和结论，突出学位论文中的创造性成果和新见解，语言力求精练。硕士学位论文中文摘要字数为500字左右。

摘要应具有独立性和完整性，能够在不阅读全文的情况下，使读者清晰了解研究问题、研究方法及主要结论。摘要一般包括以下内容：研究背景与问题提出、研究目标与研究内容、研究方法与技术路线、主要研究成果与创新点、结论与应用价值。

撰写摘要时应避免使用"本文""本章"等指代性较强的表达，一般不引用文献、不出现图表和公式。关键词应尽量选用学科领域内通用、规范的术语。

\par\vspace{\baselineskip}
\cnkeywords{硕士学位论文；论文模板；格式规范；图表使用；参考文献}
\end{cnabstract}

\cleardoublepage % 结束当前页

%%%--- "英文摘要" --- %%%%%%%%%%%%%%%%%

\begin{flushleft}
{
\renewcommand{\baselinestretch}{1}
\zihao{4}\textsf{\textbf{Title: \the\ETitle}\\
    \vspace*{14pt}
    \textbf{Major: \the\Emajor}\\
    \vspace*{14pt}
    \textbf{Name: \the\EStudentAuthor}\hfill \textbf{Signature:} \thickuline{\quad\quad\quad\quad\quad\quad}\\
    \vspace*{14pt}
    \textbf{Supervisor: \the\EsupervisorTitle~~\the\Esupervisor}\hfill \textbf{Signature:} \thickuline{\quad\quad\quad\quad\quad\quad}\\
}}
\end{flushleft}

\begin{enabstract}

The master's dissertation abstract should provide a concise summary of the research objectives, methodology, findings, and conclusions, highlighting the innovative contributions of the work. It should be approximately 500 Chinese characters in length (for the Chinese abstract).

The abstract must be self-contained and complete, enabling readers to understand the research problem, methods, and main conclusions without reading the full text. It should generally include the following elements: research background and problem statement, research objectives and content, research methods and technical approach, key findings and innovations, and conclusions with their practical or theoretical value.

When writing the abstract, avoid highly referential expressions such as "this paper" or "this chapter." Generally, do not cite references, include figures or tables, or present mathematical formulas. Keywords should prioritize standardized terminology within the discipline and be listed in order of importance.

\par\vspace{\baselineskip}
\enkeywords{Master's Dissertation; Dissertation Template; Formatting Guidelines; Use of Figures and Tables; References}
\end{enabstract}

\cleardoublepage % 结束当前页
